% Provenance (never printed): see sources/outline_sources.md and sources/book1/BRIEFS.md (g8/06).
\chapter{Project: A Data Explorer}\label{ch:g8:project-data-explorer}

The librarian of Brenmoor has two tables, the books and the loans of 2025,
and many questions. This chapter builds, in five versions, a \omterm{def:g4:programs-and-operating-system:program}{program} that
answers them. Each version runs, each is tested, and each adds one thing to
the one before. Here is what the finished \omterm{def:g4:programs-and-operating-system:program}{program} must do:

\begin{quote}
\itshape The explorer loads the library's books and loans and offers a \omterm{def:g2:using-a-computer:menu}{menu}.
It can count the loans of each genre, \omterm{def:g7:lists:list}{list} the ten most borrowed books, \omterm{def:g7:lists:list}{list}
the books never borrowed, and count the loans still out on a date that the
user types. It answers each question, and saves the last answer to a report
\omterm{def:g3:files-and-folders:file}{file}.
\end{quote}

\begin{recall}
A \omterm{def:g8:files-and-data-tables:table}{CSV file} keeps a \omterm{def:g8:files-and-data-tables:table}{data table} as text, and \code{csv.DictReader} turns it into
a \omterm{def:g7:lists:list}{list} of records (\cref{ch:g8:files-and-data-tables}). A \omterm{def:g8:dictionaries:dictionary}{dictionary} gives
the \omterm{def:g6:variables-and-input:value}{value} of a key, and counts things with \code{get}
(\cref{ch:g8:dictionaries}). \code{sorted} returns a \omterm{def:g8:searching-and-sorting:sorted}{sorted} copy of a \omterm{def:g7:lists:list}{list}
(\cref{ch:g8:searching-and-sorting}). A project grows in versions, each one
tested (\cref{ch:g7:project-quiz-game}).
\end{recall}

\section{What must it do?}

The paragraph above, turned into checks small enough to test:

\begin{enumerate}
  \item A \omterm{def:g2:using-a-computer:menu}{menu} shows the four questions and a choice to quit; any other
    choice gets a polite message, and the \omterm{def:g2:using-a-computer:menu}{menu} goes on until the user quits.
  \item The two tables are loaded, and the \omterm{def:g4:programs-and-operating-system:program}{program} checks that it has as many
    records as the \omterm{def:g3:files-and-folders:file}{files} have lines after their \omterm{def:g8:files-and-data-tables:table}{header rows}.
  \item The loans of each genre add up to the number of loans.
  \item The top ten is \omterm{def:g8:searching-and-sorting:sorted}{sorted} from the most borrowed; two books with as many
    loans are in the order of their titles.
  \item A date that is not written \code{YYYY-MM-DD} is refused with a
    message.
  \item After each answer, \code{report.txt} holds that answer.
\end{enumerate}

\begin{omfigure}
\begin{tikzpicture}[font=\scriptsize, node distance=0.85cm,
    box/.style={draw=black!60, rounded corners=2pt, align=center, minimum height=0.9cm, text width=1.6cm}]
  \node[box, fill=ocYellow!25] (csv) {\code{books.csv}\\ \code{loans.csv}};
  \node[box, fill=ocBlue!12, right=of csv] (rec) {lists of\\ records};
  \node[box, fill=ocBlue!12, right=of rec] (dic) {dictionaries\\ of counts};
  \node[box, fill=ocMint!25, right=of dic] (ans) {answer\\ lines};
  \node[box, fill=ocYellow!25, right=of ans] (rep) {\code{report.txt}};
  \draw[-{Stealth[length=2mm]}, thick] (csv) -- node[above, font=\tiny] {load} (rec);
  \draw[-{Stealth[length=2mm]}, thick] (rec) -- node[above, font=\tiny] {count} (dic);
  \draw[-{Stealth[length=2mm]}, thick] (dic) -- node[above, font=\tiny] {sort} (ans);
  \draw[-{Stealth[length=2mm]}, thick] (ans) -- node[above, font=\tiny] {save} (rep);
  \node[below=0.3cm of ans, font=\tiny] {printed on the screen too};
\end{tikzpicture}

{\small The path of the data through the explorer: from two \omterm{def:g3:files-and-folders:file}{files} to an
answer on the screen and in a report \omterm{def:g3:files-and-folders:file}{file}.}
\end{omfigure}

\section{The menu and its loop}

Version 1 has no data yet: only the \omterm{def:g2:using-a-computer:menu}{menu}, the \omterm{def:g4:algorithms-repeat-and-decide:loop}{loop} that reads the choices,
and a way out.

\omcode{code/grade-8/06-project-data-explorer/listings/explorer_v1.py}{1}{21}{Version
1: the \omterm{def:g2:using-a-computer:menu}{menu}.}\label{lst:g8:project-data-explorer:v1}

\omoutput{code/grade-8/06-project-data-explorer/outputs/explorer_v1.txt}{A session:
a question, a wrong choice, then quit.}\label{out:g8:project-data-explorer:v1}

The three quotes \code{"""} open a \omterm{def:g7:strings:string}{string} that goes on over several lines.
The \omterm{def:g4:algorithms-repeat-and-decide:loop}{loop} goes on while the choice is not \code{"0"}; the test
\code{choice in ["1", "2", "3", "4"]} checks that the choice is one of the
\omterm{def:g7:lists:list}{list}.

\begin{omfigure}
\begin{tikzpicture}[font=\scriptsize, node distance=0.45cm,
    st/.style={draw=black!60, rounded corners=2pt, fill=ocBlue!10, align=center, minimum height=0.6cm},
    dc/.style={draw=black!60, diamond, aspect=2.4, fill=ocYellow!25, inner sep=1pt, align=center}]
  \node[st] (menu) {print the menu};
  \node[st, below=of menu] (read) {read a choice};
  \node[dc, below=of read] (zero) {\code{"0"}?};
  \node[dc, below=0.6cm of zero] (ok) {1 to 4?};
  \node[st, right=1.4cm of ok] (ans) {answer it};
  \node[st, left=1.4cm of ok] (msg) {``type a number\\ from 0 to 4''};
  \node[st, right=1.4cm of zero] (bye) {goodbye};
  \draw[-{Stealth[length=1.5mm]}] (menu) -- (read);
  \draw[-{Stealth[length=1.5mm]}] (read) -- (zero);
  \draw[-{Stealth[length=1.5mm]}] (zero) -- node[above] {yes} (bye);
  \draw[-{Stealth[length=1.5mm]}] (zero) -- node[right] {no} (ok);
  \draw[-{Stealth[length=1.5mm]}] (ok) -- node[above] {yes} (ans);
  \draw[-{Stealth[length=1.5mm]}] (ok) -- node[above] {no} (msg);
  \draw[-{Stealth[length=1.5mm]}] (ans.south) -- ++(0,-0.4) -| ($(msg.south)+(-1.4,-0.4)$) |- (read.west);
  \draw (msg.south) -- ++(0,-0.4);
\end{tikzpicture}

{\small The \omterm{def:g2:using-a-computer:menu}{menu} \omterm{def:g4:algorithms-repeat-and-decide:loop}{loop}: read a choice, stop on \code{0}, answer a question,
or ask again.}
\end{omfigure}

\section{Loading the two tables}

Version 2 loads the tables before the \omterm{def:g2:using-a-computer:menu}{menu}. Two new \omterm{def:g7:functions:function}{functions} do it:

\omcode{code/grade-8/06-project-data-explorer/listings/explorer_v2.py}{10}{27}{Version
2: loading and checking the tables (the start of
\code{main}).}\label{lst:g8:project-data-explorer:v2}

\Needspace*{12\baselineskip}
\omoutput{code/grade-8/06-project-data-explorer/outputs/explorer_v2.txt}{80 books
and 400 loans, then quit at once.}\label{out:g8:project-data-explorer:v2}

\code{assert} followed by a \omterm{def:g4:algorithms-repeat-and-decide:condition}{condition} stops the \omterm{def:g4:programs-and-operating-system:program}{program} with an error if the
\omterm{def:g4:algorithms-repeat-and-decide:condition}{condition} is false, and does nothing otherwise: a check written into the
\omterm{def:g4:programs-and-operating-system:program}{program}. Here it compares the number of records with the number of lines of
the \omterm{def:g3:files-and-folders:file}{file}, so the 80 and the 400 are read from the data, never typed by hand.
If someone adds a book to the \omterm{def:g3:files-and-folders:file}{file}, the \omterm{def:g4:programs-and-operating-system:program}{program} still agrees with it.

\section{Filtering and counting}

Each loan names its book only by a number, \code{book_id}. To count the loans
of each genre, the \omterm{def:g4:programs-and-operating-system:program}{program} looks up the book of each loan. A \omterm{def:g8:dictionaries:dictionary}{dictionary}
from \code{book_id} to the book's record makes each look-up a single step,
instead of a search through the 80 books.

\begin{omfigure}
\begin{tikzpicture}[x=1.3cm, y=0.5cm, font=\scriptsize]
  \node[anchor=south west, font=\footnotesize\bfseries] at (0,1.1) {loans};
  \foreach \l/\b [count=\i from 1] in {1/59, 2/73, 3/28, 4/54} {
    \draw[black!60, fill=white] (0,-\i) rectangle ++(0.8,1); \node at (0.4,-\i+0.5) {\l};
    \draw[black!60, fill=ocYellow!30] (0.8,-\i) rectangle ++(0.9,1); \node at (1.25,-\i+0.5) {\b};
  }
  \node[font=\tiny] at (0.4,0.6) {loan\_id};
  \node[font=\tiny] at (1.25,0.6) {book\_id};
  \node[anchor=south west, font=\footnotesize\bfseries] at (5,1.1) {books};
  \foreach \b/\g [count=\i from 1] in {28/comics, 54/sport, 59/mystery, 73/adventure} {
    \draw[black!60, fill=ocYellow!30] (5,-\i) rectangle ++(0.9,1); \node at (5.45,-\i+0.5) {\b};
    \draw[black!60, fill=white] (5.9,-\i) rectangle ++(1.6,1); \node at (6.7,-\i+0.5) {\g};
  }
  \node[font=\tiny] at (5.45,0.6) {book\_id};
  \node[font=\tiny] at (6.7,0.6) {genre};
  \foreach \a/\b in {1/3, 2/4, 3/1, 4/2} {
    \draw[-{Stealth[length=1.5mm]}, ocBlue] (1.7,-\a+0.5) -- (5,-\b+0.5);
  }
\end{tikzpicture}

{\small Looking up the book of each loan: the \code{book_id} of a loan leads
to the row of its book, and so to its genre.}
\end{omfigure}

\omcode{code/grade-8/06-project-data-explorer/listings/explorer_v3.py}{22}{49}{Version
3: looking up books, loans per genre, books never
borrowed.}\label{lst:g8:project-data-explorer:v3}

The \omterm{def:g2:using-a-computer:menu}{menu} \omterm{def:g4:algorithms-repeat-and-decide:loop}{loop} now calls a \omterm{def:g7:functions:function}{function} \code{answer}, which builds the lines of an
answer:

\omcode{code/grade-8/06-project-data-explorer/listings/explorer_v3.py}{52}{67}{Version
3: the answers as lines.}\label{lst:g8:project-data-explorer:v3-answer}

\omoutput[1-2,8-20]{code/grade-8/06-project-data-explorer/outputs/explorer_v3.txt}{Questions
1 and 3.}\label{out:g8:project-data-explorer:v3}

\code{answer} returns its lines instead of printing them: the last version
will both print them and save them. Question 3 has an empty answer: every
book of the library went out at least once in 2025. Saying so is better than
printing nothing, which would look like a \omterm{def:g7:finding-and-fixing-bugs:bug}{bug}.

\begin{omfigure}
\begin{tikzpicture}
  \begin{axis}[width=10cm, height=4.8cm, ybar, bar width=0.5cm, ymin=0,
      symbolic x coords={adventure,comics,history,mystery,nature,poetry,science,sport},
      xtick=data, x tick label style={font=\scriptsize, rotate=30, anchor=north east},
      nodes near coords, nodes near coords style={font=\scriptsize},
      ylabel={loans}, label style={font=\footnotesize}, tick label style={font=\footnotesize},
      axis x line*=bottom, axis y line*=left, enlarge x limits=0.08]
    \addplot[fill=ocBlue!60, draw=ocBlue] table[x=genre, y=loans] {code/grade-8/06-project-data-explorer/figures/out/genres.dat};
  \end{axis}
\end{tikzpicture}

{\small The loans of 2025 by genre, counted by the explorer's own \omterm{def:g7:functions:function}{function}
\code{loans_by_genre}.}
\end{omfigure}

\section{Sorting and the top ten}

The ten most borrowed books need the books \omterm{def:g8:searching-and-sorting:sorted}{sorted} by their number of loans,
from the most. \code{sorted} can take a \code{key}: a \omterm{def:g7:functions:function}{function} that it calls
on each \omterm{def:g7:lists:list}{item}, and whose results it compares instead of the \omterm{def:g7:lists:list}{items}. Here the
key of a row is the \omterm{def:g7:lists:list}{list} \code{[-loans, title]}. \omterm{def:g7:lists:list}{Lists} are compared \omterm{def:g7:lists:list}{item} by
\omterm{def:g7:lists:list}{item}, like the \omterm{def:g7:strings:string}{characters} of a text (\cref{ch:g8:files-and-data-tables}): the
first \omterm{def:g7:lists:list}{items} decide, and only when they are equal do the second ones. The minus
sign puts the largest numbers of loans first; the title breaks the ties.

\omcode{code/grade-8/06-project-data-explorer/listings/explorer_v4.py}{52}{65}{Version
4: the sort key and the top ten.}\label{lst:g8:project-data-explorer:v4}

\omoutput[1,8-20]{code/grade-8/06-project-data-explorer/outputs/explorer_v4.txt}{The
top ten.}\label{out:g8:project-data-explorer:v4}

The key is a named \omterm{def:g7:functions:function}{function}, given to \code{sorted} without brackets:
\code{key=most_borrowed_first} hands over the \omterm{def:g7:functions:function}{function} itself, for
\code{sorted} to call; \code{most_borrowed_first()} would call it at once.
\code{[:10]} keeps the first ten rows of the \omterm{def:g8:searching-and-sorting:sorted}{sorted list}. Four books have 9
loans: they come in the alphabetical order of their titles.

\section{Dates and the report}

The last version adds question 4 and the report. Dates in the loans table are
written \code{YYYY-MM-DD}, so they can be compared as text
(\cref{ch:g8:files-and-data-tables}); a date typed by the user is checked
first, so that it is written the same way.

\omcode{code/grade-8/06-project-data-explorer/listings/explorer_v5.py}{69}{90}{Version
5: checking a date, the loans still out, and the
report.}\label{lst:g8:project-data-explorer:v5}

\omcode{code/grade-8/06-project-data-explorer/listings/explorer_v5.py}{121}{145}{Version
5: the \omterm{def:g2:using-a-computer:menu}{menu} \omterm{def:g4:algorithms-repeat-and-decide:loop}{loop} now saves each answer.}\label{lst:g8:project-data-explorer:v5-main}

\omoutput[1-2,8-22]{code/grade-8/06-project-data-explorer/outputs/explorer_v5.txt}{Two
dates and one refused; the report keeps the last answer.}\label{out:g8:project-data-explorer:v5}

A loan is still out on a date if it went out on or before that date and had
not come back by then: either it never came back, or it came back later. On 31
December, 24 loans are out; 22 of them never came back, and 2 came back in
2026. The \omterm{def:g4:programs-and-operating-system:program}{program} checks the date before using it, and the \omterm{def:g6:variables-and-input:value}{variable}
\code{saved} makes sure that \code{report.txt} is only read once something has
been written to it.

The tests of the chapter check each \omterm{def:g7:functions:function}{function} on the tables: that the genre
counts add up to 400, that the top ten is in order, that each version's
session prints what it should, and that the \omterm{def:g4:programs-and-operating-system:program}{program}'s copies of the two
tables are exactly the town's \omterm{def:g3:files-and-folders:file}{files}.

\section{Exercises}

\begin{exercise}[$\star$]\label{exo:g8:project-data-explorer:1}
Version 1 is given the choices \code{5}, \code{1}, then \code{0}. What does it
print after the \omterm{def:g2:using-a-computer:menu}{menu}?
\end{exercise}

\begin{exercise}[$\star$]\label{exo:g8:project-data-explorer:2}
A toy library has three books, Ash (\code{book_id} 1), Birch (2) and Cedar
(3), and five loans, of the books 2, 3, 2, 1 and 3. What are the first three
rows of \code{top_ten}?
\end{exercise}

\begin{exercise}[$\star$]\label{exo:g8:project-data-explorer:3}
Write a \omterm{def:g7:functions:function}{function} \code{books_by(books, author)} that returns the titles of
the books written by one author.
\end{exercise}

\begin{exercise}[$\star$]\label{exo:g8:project-data-explorer:4}
This version of \code{top_ten} prints The Clockwork Engine in first and second
place. Why? Fix it.

\omcode{code/grade-8/06-project-data-explorer/exercises/top_ten_bug.py}{5}{16}{A top
ten that repeats books.}\label{lst:g8:project-data-explorer:top-bug}
\end{exercise}

\begin{exercise}[$\star\star$]\label{exo:g8:project-data-explorer:5}
Write a \omterm{def:g7:functions:function}{function} \code{busiest_month(loans)} that returns the month in which
the most loans went out, and their number.
\end{exercise}

\begin{exercise}[$\star\star$]\label{exo:g8:project-data-explorer:6}
Without \code{is_date}, version 5 answers \code{22 loans still out on
30/6/2025}. Explain this number, by comparing \code{"30/6/2025"} with the
dates of the table as text.
\end{exercise}

\begin{exercise}[$\star\star$]\label{exo:g8:project-data-explorer:7}
The town also has a members table, \code{members.csv}, with a
\code{member_id} column. Write a \omterm{def:g7:functions:function}{function} \code{never_borrowers(members,
loans)} that returns the members who never borrowed a book.
\end{exercise}

\begin{exercise}[$\star\star$]\label{exo:g8:project-data-explorer:8}
For 80 books and 400 loans, how many times does \code{loans_by_genre} read
or write the dictionaries \code{finder} and \code{counts}? Count the steps of
\code{book_finder} too.
\end{exercise}

\begin{exercise}[$\star\star$]\label{exo:g8:project-data-explorer:9}
Write a \omterm{def:g7:functions:function}{function} \code{save_csv(rows, filename)} that saves the top ten as a
\omterm{def:g8:files-and-data-tables:table}{CSV file}, with the \omterm{def:g8:files-and-data-tables:table}{header row} \code{title,loans}.
\end{exercise}

\begin{exercise}[$\star\star\star$]\label{exo:g8:project-data-explorer:10}
The module \code{datetime} of Python knows the calendar: after
\code{from datetime import date}, the call
\code{date.fromisoformat("2025-01-03")} turns a text into a date, and the
difference of two dates has a number of days, \code{.days}. Write a \omterm{def:g7:functions:function}{function}
\code{average_length(loans)} that returns the average number of days a book
stays out, over the loans that came back.
\end{exercise}

\begin{exercise}[$\star\star\star$]\label{exo:g8:project-data-explorer:11}
Explain why \code{top_ten} gives the same ten rows, in the same order,
whatever the order of the lines of \code{loans.csv}.
\end{exercise}

\begin{exercise}[$\star\star\star$]\label{exo:g8:project-data-explorer:12}
Write a \omterm{def:g7:functions:function}{function} \code{best_pair(loans)} that finds the pair of books
borrowed by the largest number of the same members, and that number.
\end{exercise}
